% ============================================================================
%  Task_W02_Answers.tex -- Fillable answer sheet for Task_W02: Backtracking
%  Uses coa-accessible-student.sty for a consistent look with the
%  assignment itself, without the AI-development disclaimer.
%
%  Students: see the "Before You Start" section in the compiled PDF (right
%  after the Name line below) for how to edit and compile this file.
% ============================================================================
\documentclass[11pt]{article}
\usepackage{coa-accessible-student}

% Same plain pseudocode listing style used in the assignment itself.
\lstdefinestyle{pseudo}{
  basicstyle=\ttfamily\small\color{coaBodyText},
  frame=single,
  rulecolor=\color{coaNavy},
  breaklines=true,
  showstringspaces=false,
  columns=fullflexible,
  tabsize=2,
  numbers=none
}
\lstset{style=pseudo}

\title{Task W02: Backtracking --- Answer Sheet}
\hypersetup{pdftitle={Task W02: Backtracking -- Answer Sheet}, pdflang={en-US}}

\begin{document}

{\centering
{\Large\bfseries\color{coaNavy} Task W02: Backtracking --- Answer Sheet\par}
\vspace{0.6em}
}
\noindent Name: \underline{\hspace{3.2in}}
\vspace{1em}

\begin{center}
\rule{0.95\linewidth}{0.6pt}
\end{center}
\subsection*{Before You Start}
\noindent
This is a LaTeX source file, not a finished document; what you are reading
right now is what it looks like once compiled. To edit it and produce your
own PDF, you will need a LaTeX environment such as Overleaf, TeXworks, or a
local TeX installation; if you need help setting it up in Linux, please let
me know.

\begin{enumerate}[leftmargin=1.4em]
  \item Keep \texttt{Task\_W02\_Answers.tex} and
        \texttt{coa-accessible-student.sty} in the same folder. This file
        will not compile without the style file present.
  \item Open \texttt{Task\_W02\_Answers.tex} in a text or code editor, not a
        word processor, and replace each \textit{italicized} placeholder
        with your own answer, following the instructions given in each
        section.
  \item Compile it by running \texttt{pdflatex Task\_W02\_Answers.tex}. Run
        this command twice in a row; the page numbers at the bottom of each
        page will not resolve correctly on a single pass, and will show
        "Page 1 of ??" instead of the real total.
  \item Every time you make a change to the \texttt{.tex} file, save it and
        run \texttt{pdflatex} twice again to see your latest edits in the
        PDF you submit.
  \item Submit both files: your edited \texttt{Task\_W02\_Answers.tex}
        source and the \texttt{Task\_W02\_Answers.pdf} it produces.
\end{enumerate}
\begin{center}
\rule{0.95\linewidth}{0.6pt}
\end{center}
\vspace{0.6em}

\section{Problem 1: A Smaller Board}

\subsection{Board}
Fill in the board below with a Q in each square where a queen ends up.
Columns are 1 through 4 left to right; rows are 1 through 4 top to bottom.

The table has one line of source for each row of the board. Find the line
that starts with the row number you want, then count across its blank
cells the same way you count across the column headers above the table,
and type a capital Q into the correct one. Leave every other cell in that
line blank. Do not add, remove, or reorder any \verb|&| symbols or the
\verb|\\| at the end of a line, since that changes the shape of the table.

For example, to place a queen in row 2, column 3, change this line:
\begin{lstlisting}
2 & & & & \\
\end{lstlisting}
to this:
\begin{lstlisting}
2 & & & Q & \\
\end{lstlisting}

\begin{center}
\renewcommand{\arraystretch}{1.5}
\begin{tabular}{c|c|c|c|c}
 & 1 & 2 & 3 & 4 \\
\hline
1 & & & & \\
\hline
2 & & & & \\
\hline
3 & & & & \\
\hline
4 & & & & \\
\end{tabular}
\end{center}

\subsection{Backtrack Notes}
One line for each column where your first-tried row did not work.

\begin{itemize}[leftmargin=1.4em]
  \item \textit{Column \_\_, tried row \_\_ first, failed because \ldots}
\end{itemize}

\section{Problem 2: Row by Row}

\subsection{Your Rewritten Pseudocode}
If you used the AI-for-formatting exception described in the assignment,
submit the AI conversation as its own file alongside this answer sheet.

\begin{lstlisting}
% Write your rewritten pseudocode here.
\end{lstlisting}

\subsection{Board}
Fill in this table the same way you filled in the one for Problem 1, using
a capital Q in the correct cell of the correct row's line. This time the
queens come from your rewritten row-by-row version, starting from row 0.

\begin{center}
\renewcommand{\arraystretch}{1.5}
\begin{tabular}{c|c|c|c|c}
 & 1 & 2 & 3 & 4 \\
\hline
1 & & & & \\
\hline
2 & & & & \\
\hline
3 & & & & \\
\hline
4 & & & & \\
\end{tabular}
\end{center}

\subsection{Backtrack Notes}
One line for each row where your first-tried column did not work.

\begin{itemize}[leftmargin=1.4em]
  \item \textit{Row \_\_, tried column \_\_ first, failed because \ldots}
\end{itemize}

\section{Problem 3: Watch and Answer}

\subsection{Question 1}
\textit{In your own words, explain what backtracking means for a knight's
tour search, using the words "dead end" and "undo."}

\vspace{2em}

\subsection{Question 2}
\textit{Would a corner start need more backtracking than a central start,
less, or is it impossible to say without trying? Explain your reasoning.}

\vspace{2em}

\subsection{Question 3}
\textit{Could changing the fixed move order matter here, the way it did in
Problem 2? Explain your reasoning.}

\vspace{2em}

\subsection{Question 4}
\textit{Could a backtracking search ever find a complete tour on a three by
three board? Explain your reasoning.}

\vspace{2em}

\section{Reflection Questions}

\subsection{Question 1}
\textit{Compare your Problem 1 board to your Problem 2 board, along with
your backtrack notes for each.}

\vspace{2em}

\subsection{Question 2}
\textit{Compare your Problem 1 backtrack notes to the Knight's Tour worked
example's 31 calls and 19 backtracks.}

\vspace{2em}

\subsection{Question 3}
\textit{Look back at your answer to Problem 3, question 2. Does that
reasoning match what you saw across your Eight-Queens work?}

\vspace{2em}

\subsection{Question 4}
\textit{Would changing the row order within each column for Problem 1
behave more like Problem 2, or more like the Knight's Tour example?
Explain your reasoning.}

\section{Bugs or Issues You Found}
If you found something in the assignment that needs to be debugged, for
example an unclear problem statement or an example that does not behave
the way it is described, describe it here. Leave this blank if nothing
came up.

\vspace{3em}

\end{document}
